\documentclass[project-ja]{wbsmba}

\title{プロジェクト研究論文向けLaTeXテンプレートの試作}
\subtitle{WBS提出様式への実務対応}
\author{早稲田 光}
\studentid{57XXXXXX}
\zemi{事業創造研究ゼミ}
\completiondate{2026年3月修了}
\chiefexaminer{大隈 薫 教授}
\deputyexaminer{坪内 翼 教授}

\begin{document}
\maketitle

\begin{abstractpage}
本サンプルは、プロジェクト研究論文の提出体裁を LaTeX で扱えるようにするための試作である。
専門職学位論文と多くの要素を共有しつつ、審査教員表記などの差分を吸収できるよう設計した。
\end{abstractpage}

\wbsfrontmatter
\tableofcontents
\wbsmainmatter

\chapter{はじめに}
\section{問題設定}
プロジェクト研究論文では、実務上の課題を具体的に定義し、分析枠組みと提案の整合性を示すことが重要である。
その際、戦略論の入門書や体系書が示す基本的な論点整理は、問題設定の粒度を整えるうえで有用である
\cite{asaba2010tsukamu,ishii1996strategy}。

\chapter{分析}
\section{現状把握}
既存テンプレートの外形を踏まえつつ、LaTeX の標準機能で継続利用しやすい構成を採用する。

\section{提案}
見た目の厳密一致よりも、教員が確認しやすい体裁と、執筆時の再現性を優先する。

\chapter{結論}
このクラスは、WBS の日本語プロジェクト研究論文にもそのまま流用できる。

\bibliographystyle{wbsapalike}
\bibliography{sample-project-ja}

\appendix
\chapter{補足資料}
追加の表や質問票を付録として置く。

\end{document}
